\documentclass[11pt,a4paper]{article}
\usepackage[margin=2.5cm]{geometry}
\usepackage{amsmath,amssymb}
\usepackage{booktabs}
\usepackage{array}
\usepackage{xcolor}
\usepackage{tcolorbox}
\tcbuselibrary{skins,breakable}
\usepackage{hyperref}
\usepackage[numbers,sort&compress]{natbib}
\bibliographystyle{unsrtnat}
\usepackage{enumitem}
\hypersetup{colorlinks=true,linkcolor=blue!60!black,urlcolor=blue!60!black}
\setlength{\emergencystretch}{2em}

\definecolor{boxblue}{RGB}{220,235,255}
\definecolor{boxorange}{RGB}{255,240,210}
\definecolor{boxgreen}{RGB}{220,245,220}
\definecolor{boxyellow}{RGB}{255,252,220}
\definecolor{boxred}{RGB}{255,230,225}

\newcommand{\term}[1]{\textbf{#1}}
\newcommand{\lean}[1]{\texttt{#1}}
\newcommand{\Neff}{N_{\mathrm{eff}}}
\newcommand{\Nc}{N_c}
\newcommand{\CatAL}{\textsf{CatAL}}
\newcommand{\CatA}{\textsf{CatA}}
\newcommand{\CatAD}{\textsf{CatAD}}
\newcommand{\PHI}{\varphi}

\title{\textbf{From the Arithmetic Substrate to Particle Masses}\\[6pt]
\large The GTE Cascade and the Nine Charged-Fermion Masses: A Step-by-Step Guide\\[4pt]\normalsize\textit{UGP Physics Tutorial Series}
\\[2pt]\small\href{https://doi.org/\TutorialDOI}{doi.org/\TutorialDOI}}
\author{Nova Spivack\\[4pt]
\normalsize
\href{https://novaspivack.com/research}{novaspivack.com/research}\\[2pt]
\small Full monograph (P48):
\href{https://doi.org/10.5281/zenodo.20560550}{doi.org/10.5281/zenodo.20560550}
}
\newcommand{\TutorialDOI}{10.5281/zenodo.20657913}
\date{2026}

\begin{document}
\setcounter{tocdepth}{2}
\maketitle
\tableofcontents
\newpage

\begin{tcolorbox}[colback=blue!6,colframe=blue!40!black,
  title=\textbf{UGP Physics Tutorial Series}]
\textbf{Prerequisites (read first):} \href{https://doi.org/10.5281/zenodo.20657889}{\emph{The GTE Polynomial: A Step-by-Step Cheat Sheet}}\\[4pt]
\textbf{See also:} \href{https://doi.org/10.5281/zenodo.20657898}{\emph{How the Standard Model Quantum Numbers Emerge from $\mathbb{Z}_7\times\mathbb{Z}_3$}} $\cdot$
\href{https://doi.org/10.5281/zenodo.20657876}{\emph{How MDL Selects the 19-Bit Polynomial}}
\end{tcolorbox}

\begin{tcolorbox}[colback=boxgreen,colframe=green!50!black,
  title=\textbf{Claim-Strength Taxonomy}]
Throughout this tutorial, results are labeled with a confidence grade:
\begin{itemize}[itemsep=2pt]
  \item \textbf{$\mathbf{A}_{\rm Lean}$ (CatAL)} --- Machine-certified in Lean~4 with zero \texttt{sorry} and zero custom axioms.  Independently verifiable by anyone with Lean installed.
  \item \textbf{A/D (CatAD)} --- Analytically derived with full derivation, computationally confirmed.  Not yet machine-certified.
  \item \textbf{A (CatA)} --- Analytically derived; derivation complete but not independently machine-certified.
  \item \textbf{A/D (conditional)} --- Derived under a named, stated premise; the premise is itself an open problem or a CatB assumption.
  \item \textbf{B (CatB)} --- A stated working hypothesis or structural premise; not yet derived.
\end{itemize}
\end{tcolorbox}


%─────────────────────────────────────────────────────────────────────────────
\section{What Is the Universal Generative Principle?}
\label{sec:ugp}
%─────────────────────────────────────────────────────────────────────────────

\begin{tcolorbox}[colback=boxblue,colframe=blue!40!black,title=The one-line answer]
The \term{Universal Generative Principle} (UGP) is a number-theoretic sieve that,
starting from a single integer ridge level $n$, severely restricts which starting
triples are internally self-consistent --- and identifies our universe as the
\emph{unique} minimum-description triple among all admissible universes.
\end{tcolorbox}

\subsection{The parameter problem}

The Standard Model of particle physics requires approximately 25 numerical constants
as external inputs: quark and lepton masses, gauge couplings, mixing angles, and
cosmological parameters.  From the Standard Model's perspective, these numbers are
independent: they span a 25-dimensional space of possible universes.  There is no
known reason why they take their specific values.

The UGP shows this picture is wrong.  A number-theoretic sieve applied across
\emph{all} ridge levels $n$ constrains which parameter combinations are
arithmetically self-consistent.  Most starting points fail the sieve.  The survivors
do not fill the full 25-dimensional space; they form a lower-dimensional constraint
manifold.  Our universe occupies the unique minimum-description point on that manifold.

\subsection{Ridge structure and the $n = 10$ uniqueness}

The sieve operates at each integer \term{ridge level} $n$.  At level $n$, the
\term{ridge value} is
\[
  R_n \;=\; 2^n - 16.
\]
The constant 16 is not a free parameter: $16 = 2^{1+\mathrm{ord}_7(2)}$, where
$\mathrm{ord}_7(2) = 3$ is the multiplicative order of 2 modulo 7 (the unique
integer $k$ such that $2^k \equiv 1 \pmod 7$, here $k = 3$ since $2^3 = 8 \equiv 1$).
This follows from the $\mathbb{Z}_7 \times \mathbb{Z}_3$ algebraic structure
that MDL forces on the system and is machine-certified
(\lean{gt001\_gt009\_ridge\_closed}, \CatAL, zero sorry).

The sieve scans all interior divisor pairs $(b_2,\,q_2)$ of $R_n$ (pairs with both
factors $> 15$) and filters them through two constraints:
\begin{itemize}[itemsep=2pt]
  \item \textbf{Prime-lock:} the quantity $c_1 = b_1 q_1 + 20$ must be prime, where
    $b_1 = b_2 + q_2 + 7$ and $q_1 = b_2 - 13$.
  \item \textbf{Mirror duality:} the swapped pair $(q_2, b_2)$ must also pass the
    prime-lock test.
\end{itemize}
{\sloppy
At \emph{most} ridge levels the sieve fails: no interior divisor pair satisfies both
constraints simultaneously.  The level $n = 10$ is the lex-minimal level where both
conditions can be satisfied, certified by four independent arithmetic proofs
(\lean{n10\_is\_minimal\_admissible\_ridge},
\lean{asymptotic\_sparsity\_universal}, \CatAL, zero sorry).
\par}

At $n = 10$: $R_{10} = 2^{10} - 16 = 1008$.  Running the sieve over all five
interior divisor pairs:

\begin{center}
\small
\renewcommand{\arraystretch}{1.35}
\begin{tabular}{@{}ccccc@{}}
\toprule
$(b_2,\;q_2)$ & $b_1$ & $q_1$ & $c_1 = b_1 q_1 + 20$ & Prime? \\
\midrule
$(16,\;63)$ & $86$          & $3$  & $278 = 2\times139$          & No  \\
$(18,\;56)$ & $81$          & $5$  & $425 = 5^2\times17$         & No  \\
$(21,\;48)$ & $76$          & $8$  & $628 = 4\times157$          & No  \\
$(24,\;42)$ & $\mathbf{73}$ & $11$ & $\mathbf{823}$ (prime)      & \textbf{Yes} \\
$(28,\;36)$ & $71$          & $15$ & $1085 = 5\times7\times31$  & No  \\
\bottomrule
\end{tabular}
\end{center}

Exactly one pair, $(24,\,42)$, passes the prime-lock test.  Its mirror swap
$(42,\,24)$ also passes: $c_1' = 73\times 29 + 20 = 2137$, also prime.

\subsection{The Lepton Seed and MDL selection}

{\sloppy
The prime-lock uniqueness produces two admissible starting triples:
$(1,\,73,\,823;\,1)$ and $(1,\,73,\,2137;\,1)$.  Two candidates, not infinitely
many.  The \term{Minimum Description Length (MDL)} principle selects the shorter
description: among all admissible triples, the lexicographically minimal one is the
physical universe.
\par}

\begin{tcolorbox}[colback=boxgreen,colframe=green!50!black,title=The Lepton Seed]
The \term{Lepton Seed} is the unique MDL-minimal admissible triple at ridge level
$n = 10$:
\[
  (1,\;\;73,\;\;823;\;\;1).
\]
Components:
\begin{itemize}[itemsep=1pt]
  \item $a = 1$: parity-phase index (encodes chirality via the M\"{o}bius function).
  \item $b = 73$: ladder index (the electron's N-value; will become the mass address).
  \item $c = 823$: branch capacity (the prime selected by the prime-lock test).
  \item $g = 1$: generation index (this is the first generation).
\end{itemize}
Machine-certified: \lean{rsuc\_theorem}, \lean{mdl\_selects\_LeptonSeed}
(\CatAL, zero sorry; P01~\cite{SpivackSM}).
\end{tcolorbox}

The number 73 is not arbitrary.  It equals $\Nc^4 - (\Nc^2+1)/2 - \Nc$ evaluated at
the QCD color rank $\Nc = 3$:
\[
  b_1 = 3^4 - \tfrac{3^2+1}{2} - 3 = 81 - 5 - 3 = 73.
\]
This is the \emph{Frobenius-Generation Coincidence Identity}, machine-certified by
17 zero-sorry Lean theorems: the electron's N-value is arithmetically forced by
the number of quark colors, not fitted to data.

\medskip
\noindent\textbf{What ``MDL-optimal'' means in plain English.}\enspace
MDL (Minimum Description Length) is the principle that the correct description of a
system is the \emph{shortest complete description} --- the one that encodes the most
structure in the fewest bits.  When two arithmetic constructions are both valid, the
one requiring fewer bits to specify is the physically realized one.  In this context,
$(1,\,73,\,823;\,1)$ beats $(1,\,73,\,2137;\,1)$ because 823 requires fewer bits to
specify than 2137 --- and only the MDL-selected triple leads to the correct particle
spectrum.  The selection is a theorem, not a heuristic
(\lean{mdl\_selects\_LeptonSeed}, \CatAL).

%─────────────────────────────────────────────────────────────────────────────
\section{What Is GTE?}
\label{sec:gte}
%─────────────────────────────────────────────────────────────────────────────

\begin{tcolorbox}[colback=boxorange,colframe=orange!60!black,title=Key idea]
\term{Generative Triple Evolution} (GTE)~\cite{SpivackGTECompleteFramework} is the two-step arithmetic map that
advances an integer triple one generation at a time.  Starting from the Lepton
Seed, two applications of the map $T$ produce the three lepton generations ---
electron, muon, tau --- with no free parameters.
\end{tcolorbox}

\subsection{The map $T$: odd and even steps}

The GTE map $T$ takes a triple $(a,\,b,\,c;\,g)$ and a step type
($t = 1$ for odd, $t = 2$ for even) and produces a new triple.  The
mechanism: divide the capacity $c$ by the ladder index $b$ to get
quotient $q = \lfloor c/b \rfloor$ and remainder $m = c \bmod b$.
These two numbers control the update of all three components.

\begin{tcolorbox}[colback=boxblue,colframe=blue!40!black,
  title={The GTE map $T$ (ridge $n$, color rank $N_c$)},breakable]
\textbf{Step 1.} Divide: $c \div b \;\to\; q$ (quotient), $m$ (remainder).

\medskip
\textbf{Step 2a. Parity-phase update (both steps):}
\[
  a' = m - (n + 2 - t).
\]

\textbf{Step 2b. Ladder-index update:}
\[
  b' = \begin{cases}
    b - (m + q) & \text{odd step } (t=1) \\[4pt]
    b + F_{|q_2 - q_1|} & \text{even step } (t=2)
  \end{cases}
\]
where $F_k$ is the $k$-th Fibonacci number and $q_1, q_2$ are the quotients from
steps 1 and 2 respectively.

\textbf{Step 2c. Capacity update (Mersenne saturation):}
\[
  c' = \begin{cases}
    2^n - 1 & \text{odd step: saturate to ridge Mersenne ceiling} \\[4pt]
    2^{n + 2\Nc} - 1 & \text{even step: extend by color-rank jump}
  \end{cases}
\]
No free parameter appears anywhere.
\end{tcolorbox}

\medskip
\noindent\textbf{Why ``Generative Triple Evolution''?}\enspace
The name describes the mechanism precisely.  \emph{Generative}: the map produces a
new triple at each step, generating the particle spectrum one generation at a time.
\emph{Triple}: the objects are integer triples $(a,\,b,\,c)$.
\emph{Evolution}: the two alternating steps advance the triple through a deterministic
orbit --- like a crystal growth law where each step is forced by the geometry of the
previous one.

\medskip
\noindent\textbf{The odd step contracts; the even step lifts.}\enspace
In the odd step, the ladder index $b$ \emph{decreases}: it shrinks by the sum $m+q$.
This contraction unwinds the sieve, arithmetically recovering the interior divisor
$b_2 = 42$ from which the seed was constructed.

In the even step, $b$ \emph{increases} by a Fibonacci number.  The Fibonacci lift
is not a choice: once the quotient gap between the two steps is fixed (it is always
exactly 13 at $n = 10$), the unique lift consistent with the arithmetic constraints
is $F_{13} = 233$.  That the gap is exactly 13 is itself a theorem:
$|q_2 - q_1| = 13 = F_7 = F_{\mathrm{ord}_7(2) + \Nc + 1}$, derived from
the $\mathbb{Z}_7 \times \mathbb{Z}_3$ structure of the theory
(\lean{gt012\_gap\_from\_z7\_z3}, \CatAL, zero sorry).  Fibonacci numbers enter
because the underlying algebraic structure forces them — not by assumption.

%─────────────────────────────────────────────────────────────────────────────
\section{The Worked Cascade: Three Lepton Generations}
\label{sec:cascade}
%─────────────────────────────────────────────────────────────────────────────

\begin{tcolorbox}[colback=boxblue,colframe=blue!40!black,title=Orientation]
We now follow the Lepton Seed through two applications of $T$, producing the muon
triple and the tau triple.  Every number below can be verified with pencil and paper.
The key output is the set of \term{N-values} $\{73,\,42,\,275\}$, which will become
the mass addresses for the three charged leptons.
\end{tcolorbox}

\subsection{Starting point: the Lepton Seed}

\[
  (a,\,b,\,c;\,g) \;=\; (1,\;\;73,\;\;823;\;\;1), \qquad n = 10, \quad \Nc = 3.
\]

\subsection{Step 1 (odd step, $t=1$): Electron $\to$ Muon generation}

\textbf{Divide:} $823 \div 73 = 11$ remainder $20$.  So $q_1 = 11$, $m_1 = 20$.

\textbf{Update each component:}
\begin{align*}
  b' &= b - (m_1 + q_1) = 73 - (20 + 11) = 73 - 31 = \mathbf{42}, \\[4pt]
  c' &= 2^{10} - 1 = 1024 - 1 = \mathbf{1023}, \\[4pt]
  a' &= m_1 - (n + 2 - t) = 20 - (10 + 2 - 1) = 20 - 11 = \mathbf{9}.
\end{align*}

\begin{tcolorbox}[colback=boxgreen,colframe=green!50!black,title=Result: Muon triple]
\[
  (9,\;\;42,\;\;1023;\;\;2)
\]
The new N-value is $b' = \mathbf{42}$ --- the muon's arithmetic address.
\end{tcolorbox}

\noindent\textbf{Why is each number forced?}
\begin{itemize}[itemsep=2pt]
  \item $b' = 42$ is the interior divisor $b_2 = 42$ from the original sieve: the
    odd step arithmetically unwinds the sieve construction.
  \item $c' = 1023 = 2^{10} - 1$ is the Mersenne number at ridge level $n = 10$:
    forced by the axioms.
  \item $a' = 9 = \Nc^2 = 3^2$: a machine-certified algebraic identity, not an
    assignment.
  \item $m_1 = 20$: forced by the prime-lock constraint
    ($c_1 = b_1 q_1 + 20 = 73 \times 11 + 20 = 823$).
\end{itemize}

\subsection{Step 2 (even step, $t=2$): Muon $\to$ Tau generation}

Start from $(9,\,42,\,1023;\,2)$.

\textbf{Divide:} $1023 \div 42 = 24$ remainder $15$.  So $q_2 = 24$, $m_2 = 15$.

\medskip
\noindent\textit{Why is $m_2 = 15$ forced?}  The \emph{ridge remainder lock theorem}
proves: for any divisor $b$ of $R_n = 2^n - 16$,
$(2^n - 1) \bmod b = 15$ whenever $n \geq 5$.  Since $b_2 = 42$ divides
$R_{10} = 1008$, the remainder $15$ is structurally guaranteed
(\lean{ridge\_remainder\_lock\_general}, \CatAL, zero sorry).

\medskip
\textbf{The quotient gap:} $|q_2 - q_1| = |24 - 11| = 13$.  This gap equals
$m_1 - 7 = 20 - 7 = 13$ --- a theorem, not a coincidence.

\textbf{The 13th Fibonacci number:}
\[
  1,\;1,\;2,\;3,\;5,\;8,\;13,\;21,\;34,\;55,\;89,\;144,\;\mathbf{233}.
  \qquad F_{13} = 233.
\]

\textbf{Update each component:}
\begin{align*}
  b' &= b + F_{13} = 42 + 233 = \mathbf{275}, \\[4pt]
  c' &= 2^{10 + 2 \times 3} - 1 = 2^{16} - 1 = 65536 - 1 = \mathbf{65535}, \\[4pt]
  a' &= m_2 - (n + 2 - t) = 15 - (10 + 2 - 2) = 15 - 10 = \mathbf{5}.
\end{align*}

\begin{tcolorbox}[colback=boxgreen,colframe=green!50!black,title=Result: Tau triple]
\[
  (5,\;\;275,\;\;65535;\;\;3)
\]
The new N-value is $b' = \mathbf{275}$ --- the tau's arithmetic address.
\end{tcolorbox}

\noindent\textbf{Why is each number forced?}
\begin{itemize}[itemsep=2pt]
  \item $b' = 275$: once the quotient gap is forced to 13, the Fibonacci lift
    $F_{13} = 233$ is uniquely determined, and $42 + 233 = 275$.
  \item $c' = 65535 = 2^{16} - 1$: the exponent jump $2\Nc = 6$ at $\Nc = 3$
    carries the ridge from $n = 10$ to $n + 2\Nc = 16$.
  \item $a' = 5 = (\Nc^2 + 1)/2$: another structural identity, machine-certified.
    The general $a$-value formula $a' = m - (n + 2 - t)$ used at both steps is
    proved to be an \emph{orbit projection} — a consequence of the cascade
    arithmetic, not an ad hoc assignment
    (\lean{semantic\_floor\_a\_condition\_from\_cascade}, \CatAL, zero sorry).
\end{itemize}

\subsection{Summary: the three lepton generations}

\begin{tcolorbox}[colback=boxorange,colframe=orange!60!black,
  title={Lepton Cascade Result (\CatAL, \lean{update\_map\_produces\_canonical\_orbit})}]
\begin{center}
\renewcommand{\arraystretch}{1.35}
\begin{tabular}{@{}clcc@{}}
\toprule
Generation & Triple $(a,\,b,\,c;\,g)$ & $\Neff = b$ & Particle \\
\midrule
$g = 1$ (seed)      & $(1,\;\;73,\;\;\;\;823;\;1)$ & $\mathbf{73}$  & Electron \\
$g = 2$ (odd step)  & $(9,\;\;42,\;\;1023;\;2)$    & $\mathbf{42}$  & Muon     \\
$g = 3$ (even step) & $(5,\;275,\;65535;\;3)$       & $\mathbf{275}$ & Tau      \\
\bottomrule
\end{tabular}
\end{center}
The three N-values $\{73,\,42,\,275\}$ are the output of two forced arithmetic
steps from a sieve-forced seed.  They are not fitted.  Zero dials.
\end{tcolorbox}

%─────────────────────────────────────────────────────────────────────────────
\section{From N-Values to Masses: The InformationMassTransformer}
\label{sec:imt}
%─────────────────────────────────────────────────────────────────────────────

\begin{tcolorbox}[colback=boxblue,colframe=blue!40!black,title=The pipeline in one equation]
\[
  m_f \;=\; \mathcal{C}_f(a,\,b,\,c;\,g) \;\times\; E_{\rm base}(N_{\rm eff},\,g).
\]
Two factors.  Each is a deterministic function of the triple's arithmetic content.
No parameter is adjusted at prediction time.
\end{tcolorbox}

\subsection{The base energy $E_{\rm base}(N_{\rm eff},\,g)$}

The \term{base energy} $E_{\rm base}(\Neff, g)$ is an information-theoretic energy
scale assembled deterministically from the triple's N-value and generation index.
Five locked components contribute: an entropy term (proportional to the Shannon
information content of $\Neff$ bits), a holographic-radius term (from the area
encoding of the orbit), a coherence term, a phase term, and a binding term.
Each component is an explicit locked function of $\Neff$ and $g$, with no free
parameters (P01).

The physical picture: in GTE, a particle's rest mass equals the minimum energy
required to sustain a stable kink whose PSC orbit address has information content
$\Neff$ bits.  A larger address (more bits) means more structural information must
be ``carried'' by the kink, requiring more rest energy; this is why muons are heavier
than electrons ($\Neff^{\mu} = 42 < \Neff^{e} = 73$ would suggest the wrong order,
but the generation index $g$ in the base energy formula accounts for the mass hierarchy correctly).

The base energy sets the \emph{scale} at which each particle ``weighs.''  It
is the information-theoretic counterpart of a mass eigenvalue: the amount of
mass-energy that an $\Neff$-bit informational address at generation $g$ naturally
carries under the GTE encoding.

For the three lepton addresses, the engine evaluates to:

\begin{center}
\renewcommand{\arraystretch}{1.35}
\begin{tabular}{@{}llcc@{}}
\toprule
Particle & Triple & $\Neff$ & $E_{\rm base}(\Neff, g)$ (MeV) \\
\midrule
Electron & $(1,\;73,\;823;\;1)$    & $73$  & $0.4585$ \\
Muon     & $(9,\;42,\;1023;\;2)$   & $42$  & $110.952$ \\
Tau      & $(5,\;275,\;65535;\;3)$ & $275$ & $6534.094$ \\
\bottomrule
\end{tabular}
\end{center}

\subsection{The calibration factor $\mathcal{C}_f$}

The \term{calibration factor} $\mathcal{C}_f$ is a dimensionless number computed
from the triple's arithmetic content alone, encoding how the particle's specific
number-theoretic signature modifies the base scale.  It is computed in two steps.

\subsubsection*{Step A: The feature vector}

From the triple $(a, b, c;\, g)$, form the logarithmic ratio
\[
  L(a,b,c) = \ln\!\left(\frac{|b|}{|c|}\right),
\]
and the nine-component \emph{feature vector}:
\[
  \boldsymbol{\phi} = \bigl(1,\; L,\; L^2,\; g,\; g^2,\; M,\; \mu(a),\; \mu(b),\; \mu(|c|)\bigr),
\]
where $\mu(\cdot)$ is the \term{M\"{o}bius function}:
\[
  \mu(n) = \begin{cases}
    1  & n = 1 \\
    (-1)^k & n \text{ is a product of } k \text{ distinct primes} \\
    0  & n \text{ has a repeated prime factor}
  \end{cases},
\]
and $M = \mu(a)\,\mu(b)\,\mu(|c|)$ is the \term{M\"{o}bius product}.

The M\"{o}bius function captures the number-theoretic texture of each component:
primes give $-1$, squarefree composites give $\pm 1$ depending on parity, and
numbers with squared factors give $0$.  The feature vector encodes how the
triple's arithmetic character --- its logarithmic ratio, generation, and
number-theoretic texture --- maps to the calibration factor.

\medskip
\noindent\textbf{The M\"{o}bius product as a species classifier.}\enspace
The product $M = \mu(a)\,\mu(b)\,\mu(|c|)$ provides a number-theoretic fingerprint for
each GTE generation-orbit triple.
For the electron triple $(1, 73, 823)$: every component is either 1 or prime (hence
squarefree), giving $M = \mu(1)\cdot\mu(73)\cdot\mu(823) = (+1)(-1)(-1) = +1$.
For the muon triple $(9, 42, 1023)$: the component $a = 9 = 3^2$ has a repeated prime
factor, so $\mu(9) = 0$ and $M = 0$.
Similarly for the tau triple $(5, 275, 65535)$: $b = 275 = 5^2 \cdot 11$ gives $\mu(275) = 0$
and $M = 0$.
The electron triple is therefore the unique charged-lepton triple with $M = +1$, giving it a
distinct arithmetic signature from the muon and tau.

A complementary sector distinction appears in the lepton $c$-values.  All three satisfy
$c \equiv 7 \pmod{8}$: we have $823 = 102 \times 8 + 7$, $1023 = 127 \times 8 + 7$,
and $65535 = 8191 \times 8 + 7$.  This common residue class is a GTE-derivable property
of the lepton branch, distinguishing it arithmetically from the quark sector
(for example, the quark $c$-value 42 satisfies $42 \equiv 2 \pmod{8}$).
Lean cert: \lean{lepton\_c\_values\_mod8\_seven} (\CatAL)~\cite{SpivackGTECompleteFramework}.

\subsubsection*{Step B: The Universal Calibration Law}

\begin{tcolorbox}[colback=boxyellow,colframe=yellow!75!black,
  title={Universal Calibration Law (UCL)}, breakable]
\[
  \ln \mathcal{C}_f = \mathbf{k} \cdot \boldsymbol{\phi},
\]
where $\mathbf{k} = (k_0,\, k_L,\, k_{L^2},\, k_g,\, k_{g^2},\, k_M,\, k_a,\, k_b,\, k_c)$
is a \emph{single, fixed} nine-component coefficient vector, the same for every
fermion, locked before any prediction is made.  The locked vector used in
the canonical pipeline is:
\[
  \mathbf{k} = \begin{aligned}[t]
    (&-0.15487,\; 0.01970,\; 0.01357,\; 1.54480,\; -0.80925,\\
     &-0.80587,\; 0.12373,\; -1.50453,\; 1.32657).
  \end{aligned}
\]
\end{tcolorbox}

This vector has algebraic structure known as the \term{Elegant Kernel}: most
components are expressible in terms of the golden ratio
$\varphi = (1+\sqrt{5})/2$, the prime modulus 7, and $\pi$.  For instance,
$k_g \approx \varphi\cos(\pi/10)$, $k_{g^2} \approx -\varphi/2$, and
$k_0' \approx -1/(2\pi)$.  The Quarter-Lock identity
$k_M = k_{g^2} + \tfrac{1}{4}k_{L^2}$ is machine-certified in Lean~4
(\lean{quarterLockLaw}, \CatAL, zero sorry).
The engine constants also satisfy a structural identity:
$g_1 + 2w_{\rm phase} = 2^{N_{\rm fam}} - \pi/2^{n_{\rm ridge}}$, where
the Euler constant $e$ cancels exactly, leaving a combination expressible
entirely from GTE-derivable quantities (\lean{imt\_gen1\_phase\_structural\_pair}, \CatAL).

\medskip
\noindent\textbf{Why is $m_\tau$ the dimensional anchor?}\enspace
The base energy engine produces dimensionless ratios; one dimensional anchor is
needed to set the overall mass scale.  The tau mass $m_\tau$ is the unique
\emph{self-consistency condition} (SCC) fixed point of the lepton cascade: it is
the value at which the cascade's output, fed back as the potential parameter,
reproduces itself
(\lean{mphi\_equals\_tau\_mass\_scc}, \CatAL, zero sorry).
Nothing is fitted: $m_\tau$ is selected by arithmetic self-consistency, not by
comparison to data.

%─────────────────────────────────────────────────────────────────────────────
\section{Worked Lepton-Mass Example}
\label{sec:worked}
%─────────────────────────────────────────────────────────────────────────────

We now execute the pipeline for all three charged leptons, carrying full arithmetic.
All intermediate steps use the canonical locked coefficient vector above.

\subsection{Electron: $(1,\;73,\;823;\;1)$}

\textbf{M\"{o}bius values:}
$73$ is prime, so $\mu(73) = -1$.
$823$ is prime, so $\mu(823) = -1$.
$\mu(1) = 1$ (by definition).

\textbf{Logarithmic ratio:}
$L = \ln(73/823) = \ln(0.08870\ldots) = -2.4225$.  $L^2 = 5.8685$.

\textbf{M\"{o}bius product:}
$M = \mu(1)\cdot\mu(73)\cdot\mu(823) = 1\cdot(-1)\cdot(-1) = 1$.

\textbf{Feature vector} ($g=1$, $g^2=1$):
\[
  \boldsymbol{\phi} = (1,\; -2.4225,\; 5.8685,\; 1,\; 1,\; 1,\; 1,\; -1,\; -1).
\]

\textbf{Log-calibration factor:}
\begin{align*}
  \ln \mathcal{C}_e &= (-0.15487)(1) + (0.01970)(-2.4225) + (0.01357)(5.8685) \\
    &\quad + (1.54480)(1) + (-0.80925)(1) + (-0.80587)(1) \\
    &\quad + (0.12373)(1) + (-1.50453)(-1) + (1.32657)(-1) \\
    &= -0.15487 - 0.04773 + 0.07963 + 1.54480 - 0.80925 - 0.80587 \\
    &\quad + 0.12373 + 1.50453 - 1.32657 = \mathbf{0.1084}.
\end{align*}

\textbf{Calibration factor:} $\mathcal{C}_e = e^{0.1084} = \mathbf{1.1145}$.

\textbf{Mass:}
\[
  m_e = 1.1145 \times 0.4585~\text{MeV} = \mathbf{0.5110~\text{MeV}}.
  \qquad (\text{PDG: } 0.5110~\text{MeV}.)
\]

\subsection{Muon: $(9,\;42,\;1023;\;2)$}

\textbf{M\"{o}bius values:}
$9 = 3^2$ has a repeated factor, so $\mu(9) = 0$.
$42 = 2\cdot3\cdot7$ is squarefree with three prime factors, so $\mu(42) = (-1)^3 = -1$.
$1023 = 3\cdot11\cdot31$ is squarefree with three prime factors, so $\mu(1023) = -1$.

\textbf{Logarithmic ratio:}
$L = \ln(42/1023) = -3.1928$.  $L^2 = 10.194$.

\textbf{M\"{o}bius product:}
$M = \mu(9)\cdot\mu(42)\cdot\mu(1023) = 0\cdot(-1)\cdot(-1) = 0$.

\textbf{Feature vector} ($g=2$, $g^2=4$):
\[
  \boldsymbol{\phi} = (1,\; -3.1928,\; 10.194,\; 2,\; 4,\; 0,\; 0,\; -1,\; -1).
\]

\textbf{Log-calibration factor:}
\begin{align*}
  \ln \mathcal{C}_\mu &= (-0.15487)(1) + (0.01970)(-3.1928) + (0.01357)(10.194) \\
    &\quad + (1.54480)(2) + (-0.80925)(4) + (-0.80587)(0) \\
    &\quad + (0.12373)(0) + (-1.50453)(-1) + (1.32657)(-1) \\
    &= -0.15487 - 0.06292 + 0.13843 + 3.08960 - 3.23700 - 0 \\
    &\quad + 0 + 1.50453 - 1.32657 = \mathbf{-0.0488}.
\end{align*}

\textbf{Calibration factor:} $\mathcal{C}_\mu = e^{-0.0488} = \mathbf{0.9524}$.

\textbf{Mass:}
\[
  m_\mu = 0.9524 \times 110.952~\text{MeV} = \mathbf{105.67~\text{MeV}}.
  \qquad (\text{PDG: } 105.66~\text{MeV}.)
\]

\subsection{Tau: $(5,\;275,\;65535;\;3)$}

\textbf{M\"{o}bius values:}
$5$ is prime, so $\mu(5) = -1$.
$275 = 5^2\cdot11$ has a repeated factor, so $\mu(275) = 0$.
$65535 = 3\cdot5\cdot17\cdot257$ is squarefree with four prime factors,
so $\mu(65535) = (-1)^4 = +1$.

\textbf{Logarithmic ratio:}
$L = \ln(275/65535) = -5.4736$.  $L^2 = 29.960$.

\textbf{M\"{o}bius product:}
$M = \mu(5)\cdot\mu(275)\cdot\mu(65535) = (-1)\cdot0\cdot(+1) = 0$.

\textbf{Feature vector} ($g=3$, $g^2=9$):
\[
  \boldsymbol{\phi} = (1,\; -5.4736,\; 29.960,\; 3,\; 9,\; 0,\; -1,\; 0,\; 1).
\]

\textbf{Log-calibration factor:}
\begin{align*}
  \ln \mathcal{C}_\tau &= (-0.15487)(1) + (0.01970)(-5.4736) + (0.01357)(29.960) \\
    &\quad + (1.54480)(3) + (-0.80925)(9) + (-0.80587)(0) \\
    &\quad + (0.12373)(-1) + (-1.50453)(0) + (1.32657)(1) \\
    &= -0.15487 - 0.10783 + 0.40656 + 4.63440 - 7.28325 - 0 \\
    &\quad - 0.12373 + 0 + 1.32657 = \mathbf{-1.3022}.
\end{align*}

\textbf{Calibration factor:} $\mathcal{C}_\tau = e^{-1.3022} = \mathbf{0.2719}$.

\textbf{Mass:}
\[
  m_\tau \approx 0.2719 \times 6534.09~\text{MeV} \approx 1776.9~\text{MeV}.
\]
{\sloppy
The canonical pipeline, carrying full-precision $\mathbf{k}$ coefficients rather than
the five-digit rounded values displayed here, yields $m_\tau = \mathbf{1775.22}$~MeV
(PDG: $1776.86$~MeV, relative error $-0.0925\%$).  The canonical
calibration factors are $\mathcal{C}_e = 1.1145$, $\mathcal{C}_\mu = 0.9524$,
$\mathcal{C}_\tau = 0.2717$ (source: \texttt{dual\_path\_comparison.json}, P01).
\par}

\subsection{The three-lepton result}

\begin{tcolorbox}[colback=boxorange,colframe=orange!60!black,
  title={Three Charged Lepton Masses from the GTE Cascade}]
\begin{center}
\renewcommand{\arraystretch}{1.35}
\begin{tabular}{@{}lcccc@{}}
\toprule
Particle & $\Neff$ & GTE prediction (MeV) & PDG (MeV) & Rel.\ error \\
\midrule
Electron & $73$  & $0.510990$ & $0.510999$ & $-0.0018\%$ \\
Muon     & $42$  & $105.667$  & $105.658$  & $+0.0085\%$ \\
Tau      & $275$ & $1775.22$  & $1776.86$  & $-0.0925\%$ \\
\bottomrule
\end{tabular}
\end{center}
Three generations, one fixed coefficient vector, zero fitted parameters.
\end{tcolorbox}

%─────────────────────────────────────────────────────────────────────────────
\section{The Full Nine-Fermion Table}
\label{sec:ninefermion}
%─────────────────────────────────────────────────────────────────────────────

The same pipeline extends to the six quarks.  Each quark family has its own canonical
triple from the GTE cascade at $n = 10$ (the family-permuted seeds of P01), but the
\emph{map} --- base energy engine plus UCL calibration with the same coefficient
vector $\mathbf{k}$ --- is identical.  The full nine-fermion result is:

\begin{table}[htbp]
\centering
\caption{Nine charged-fermion masses: GTE theoretical-path predictions vs.\ PDG.
Zero continuous parameters are fitted at prediction time; $m_\tau$ is the
single self-consistency anchor.  Total RMS relative error: $\mathbf{0.295\%}$
(PDG~2022 targets).  Recomputed against PDG~2024, the RMS improves to
$\mathbf{0.261\%}$ (the downward top-mass revision moves the PDG target toward
the prediction).  Eight of the nine fermion errors are below $0.10\%$; the
RMS is dominated by the top quark.
Source: canonical artifact \texttt{dual\_path\_comparison.json} (P01).}
\label{tab:nine_fermions}
\renewcommand{\arraystretch}{1.12}
\begin{tabular}{@{}lrrr@{}}
\toprule
\textbf{Particle} & \textbf{GTE predicted (MeV)} & \textbf{PDG (MeV)} & \textbf{Rel.\ error (\%)} \\
\midrule
Electron    & $0.510990$       & $0.510999$       & $-0.0018$ \\
Muon        & $105.667$        & $105.658$        & $+0.0085$ \\
Tau         & $1{,}775.22$     & $1{,}776.86$     & $-0.0925$ \\
\midrule
Up quark    & $2.16097$        & $2.16000$        & $+0.0450$ \\
Down quark  & $4.67129$        & $4.67000$        & $+0.0277$ \\
Strange     & $93.4154$        & $93.4000$        & $+0.0165$ \\
Charm       & $1{,}275.53$     & $1{,}275.00$     & $+0.0414$ \\
Bottom      & $4{,}179.49$     & $4{,}180.00$     & $-0.0122$ \\
Top         & $171{,}159$      & $172{,}760$      & $-0.9265$ \\
\midrule
\multicolumn{3}{l}{\textbf{RMS (nine fermions, PDG 2022)}} & $\mathbf{0.295\%}$ \\
\multicolumn{3}{l}{\textbf{RMS (nine fermions, PDG 2024)}} & $\mathbf{0.261\%}$ \\
\bottomrule
\end{tabular}
\end{table}

\subsection{Honest grading of the result}

\begin{tcolorbox}[colback=boxred,colframe=red!50!black,title=Certification status]
\begin{itemize}[itemsep=3pt]
  \item \textbf{The cascade triples and the Lepton Seed:} \CatAL\ ---
    machine-certified in Lean~4 with zero sorry, zero custom axioms.
    The output integers $\{73, 42, 275\}$ are proved forced.
  \item \textbf{The structural pipeline (UCL + Elegant Kernel identities):}
    \CatA\ --- computationally confirmed, Quarter-Lock identity Lean-certified.
  \item \textbf{The $0.295\%$ RMS figure:} \CatAD\ --- computational (Category~A)
    plus a disclosed URC renormalization correction layer (Category~A/D).
    The nine mass values are predictions of a zero-active-parameter pipeline, not
    fits; the URC layer modifies the renormalization constant used in $E_{\rm base}$.
  \item \textbf{Six quark PDG containment intervals:} separately \CatAL\ --- each
    quark's mass is machine-certified to lie within its PDG~2024 uncertainty band
    (\lean{m\_u\_pdg\_interval} through \lean{m\_t\_pdg\_interval}).
\end{itemize}
A/D means the result passes every quantitative test available, with a disclosed
but not yet fully formalized correction step.  It is not a fit; it is a prediction
with a certified structural chain and an honest statement about which link is
not yet formally Lean-certified.
\end{tcolorbox}

\subsection{The 1000-permutation null test}

Could the $0.295\%$ agreement be a coincidence of the functional form?  The canonical
pipeline's 1000-permutation null suite applies the same locked UCL functional form to
1000 random permutations of the triple arithmetic --- randomly reassigning N-values
among particles --- and re-scores each permutation.  The canonical result's primary
$\sigma = 2.95\times10^{-5}$ falls \emph{outside} the entire shuffled-null
distribution: the minimum permutation $\sigma$ exceeds $5.6\times10^{-2}$, more
than three orders of magnitude above the canonical baseline.  The agreement is not
a forgiving ansatz; it is specific to the canonical triples.

%─────────────────────────────────────────────────────────────────────────────
\section{Why This Works}
\label{sec:why}
%─────────────────────────────────────────────────────────────────────────────

\subsection{MDL selection generates mass-like structures}

The UGP does not start by asking ``what should the electron mass be?''  It asks
``what is the shortest self-consistent arithmetic structure?''  The answer happens to
produce, as a byproduct, a sequence of integers $\{73, 42, 275\}$ whose information
content, when run through the InformationMassTransformer, matches the lepton mass
spectrum to four significant figures.  The agreement is not engineered; it is a
structural consequence of MDL minimality at $n = 10$.

The reason MDL produces mass-like structures is that mass --- in information-theoretic
terms --- is the cost of self-description.  A particle that requires more information
to specify (higher $\Neff$, more complex triple arithmetic) has higher base energy
and therefore higher mass.  The cascade's N-values encode this information cost
directly: the muon's $\Neff = 42$ is smaller than the tau's $\Neff = 275$, and
correspondingly the muon's base energy ($110.95$~MeV) is smaller than the tau's
($6534$~MeV).  The calibration factor $\mathcal{C}_f$ then corrects for the specific
number-theoretic texture of each triple (prime structure, generation index, etc.).

\subsection{Prime-lock and the mass ratios}

The prime-lock condition forces $c_1 = 823$ to be prime and $c_1' = 2137$ to be prime.
The effect propagates through the cascade: because $b_1 = 73$ is also prime, and
$b_2 = 42 = 2\cdot3\cdot7$ is deeply composite, and $b_3 = 275 = 5^2\cdot11$ has
a squared factor, the three triples have distinctly different M\"{o}bius signatures
($\mu(b) \in \{-1, -1, 0\}$).  The UCL's M\"{o}bius components pick up this texture
and adjust the calibration factors accordingly.  Without the prime-lock, the
M\"{o}bius patterns would be arbitrary; with it, they are arithmetically forced, and
the mass ratios come out right.

\subsection{Why three generations and not more}

The cascade terminates at three generations because the sieve produces exactly three
lepton triples before the arithmetic structure saturates.  There is no fourth step
that continues the pattern: the Mersenne capacity $c' = 2^{16} - 1 = 65535$ at
generation 3 is an arithmetic ceiling above which the prime-lock and mirror-duality
conditions fail.  The three-generation structure of the Standard Model is therefore
forced, not assumed
(\lean{update\_map\_produces\_canonical\_orbit}, \CatAL).

\subsection{The role of the golden ratio and $\pi$}

The Elegant Kernel's coefficients --- $k_g \approx \varphi\cos(\pi/10)$,
$k_{g^2} = -\varphi/2$, $k_0' = -1/(2\pi)$ --- emerge from GTE symmetry constraints,
not from fitting.  The golden ratio $\varphi$ appears because the competing
sub-dynamics of the cascade settle into a gearing ratio of $3/2$ (the logarithmic
ratio $L^* = -3/2\ln\varphi$, proved from the GTE dynamic-equilibrium condition).
The factor $\pi$ enters via Bekenstein--Fisher normalization of the information
content.  The specific combination $\varphi\cos(\pi/10)$ is the generation-scaling
coefficient forced by the arithmetic geometry of the orbit.

\subsection{The big picture: a single structure, two roles}

The GTE orbit plays two distinct roles in the theory:
\begin{enumerate}[itemsep=3pt]
  \item \textbf{Spectrum arm:} the N-values $\{73, 42, 275\}$ from the cascade
    triples encode the lepton mass addresses.  Fed through the IMT, they produce the
    charged lepton masses.  The quark N-values, from parallel family-permuted cascades,
    produce the six quark masses.
  \item \textbf{Dynamics arm:} the \emph{parity shadow} of the same orbit --- the
    binary parities of the fifteen canonical cascade triples --- pins the cellular
    automaton update rule that governs the dynamics of the $\Phi_{\rm MDL}$ field.
    That rule is $p(L,C,R) = C + R - CR - LCR \pmod{7}$, which restricts to
    Rule~110 on binary inputs.
\end{enumerate}
One integer orbit carries both the particle spectrum and the field dynamics.  The two
arms converge on the same 19-bit polynomial, a convergence called the UGP--$p$
triangle.  This structural economy --- one arithmetic substrate, two physically
distinct roles --- is what makes the GTE framework predictive rather than numerological.

\medskip
\noindent\textbf{The closed circle.}\enspace
Two independent derivation chains both determine the same polynomial $p$.
\textbf{Chain A} follows PSC arithmetic: $\text{PSC} \to \Nc = 3 \to b_1 = 73 \to$
Lepton Seed $\to$ orbit $\{73, 42, 275\} \to p$ (via the Direct-Interpolation Lift
Theorem).
\textbf{Chain B} follows MDL algebra: $\text{MDL} \to \mathbb{Z}_7 \times \mathbb{Z}_3
\to p(L,C,R) = C + R - CR - LCR$ (directly, from the structure of the field).
Both routes are machine-certified with zero sorry; they are not two independent
hypotheses about $p$ but two rigorous proofs that the arithmetic and the algebra are
consistent descriptions of the same object.  The constants discussed above --- the
ridge offset 16, the quotient gap 13, and the Fibonacci lift $F_{13} = 233$ --- are
features of Chain A that are simultaneously explained by Chain B as consequences
of $\mathbb{Z}_7 \times \mathbb{Z}_3$.

\begin{tcolorbox}[colback=boxblue,colframe=blue!40!black,
  title=One-paragraph summary of the mass chain]
Starting from the single prime $823$ (forced by the number-theoretic sieve at
ridge $n=10$), two forced arithmetic steps produce the integers $\{73,\,42,\,275\}$.
These integers, when supplied to the InformationMassTransformer --- a locked
information-theoretic map with no adjustable parameters at prediction time --- yield
the electron, muon, and tau masses to four significant figures.  The same pipeline
applied to family-permuted seeds gives the six quark masses, achieving $0.295\%$
RMS across all nine fermions.  The Lean~4 formalization certifies the cascade
arithmetic with zero sorry and zero custom axioms; the mass predictions are
computationally certified and partly Lean-certified.
\end{tcolorbox}


\section*{Further Reading}
\begin{tcolorbox}[colback=orange!8,colframe=orange!50!black,
  title=\textbf{Primary Sources}]
The results in this tutorial are derived in the following papers,
all available open-access on Zenodo and at
\href{https://novaspivack.com/research}{novaspivack.com/research}:
\begin{itemize}
  \item \textbf{P01 --- Standard Model Parameters from the UGP} (charged fermion masses):\\
        \href{https://doi.org/10.5281/zenodo.20168787}{doi.org/10.5281/zenodo.20168787}
  \item \textbf{P48 --- The Complete GTE Framework} (main synthesis monograph):\\
        \href{https://doi.org/10.5281/zenodo.20560550}{doi.org/10.5281/zenodo.20560550}
\end{itemize}
Lean~4 machine proofs: \href{https://github.com/novaspivack/ugp-lean}{github.com/novaspivack/ugp-lean}
\end{tcolorbox}

\bibliography{../../papers/bib/Spivack_Papers_Bibliography}
\end{document}
